\doublespacing % Do not change - required

\chapter{Synthesis and Discussion}
\label{ch:synthesis}

%%%%%%%%%%%%%%%%%%%%%%%%%%%%%%%%%%%%%%%
% IMPORTANT
\begin{spacing}{1} %THESE FOUR
\minitoc % LINES MUST APPEAR IN
\end{spacing} % EVERY
\thesisspacing % CHAPTER
% COPY THEM IN ANY NEW CHAPTER
%%%%%%%%%%%%%%%%%%%%%%%%%%%%%%%%%%%%%%%

\section{Cross-chapter synthesis}

Use this chapter to connect the four paper-based chapters into a single doctoral argument. The synthesis should make explicit what the combined body of work establishes that no single paper establishes alone.

\section{Answering the research questions}

Return to the research questions from Chapter~\ref{ch:introduction} and answer them using evidence from Chapters~\ref{ch:paper1}--\ref{ch:paper4}.

\section{Limitations}

Discuss limitations across the whole thesis. This should include limitations inherited from individual papers and broader limitations that only become clear when the work is viewed as a combined programme of research.

\section{Future work}

Identify research directions that follow from the whole thesis, not only from the final paper.
